% GENERATED FILE. Edit canonical lessons, then run npm run generate:books.
% canonical-source-hash: fnv1a64:f20ffc9de0ea19c1
% canonical-lessons: GE-C532-wohngemeinschaft, GE-C532-mitbewohner, GE-C532-hausordnung, GE-C532-hausmeister, GE-C532-treppenhaus

\chapter{Wohnen in der WG}
\label{ch:ge-wohnen-in-der-wg}
\hlchaptermodality{532}

\begin{chapteropening}
\textbf{By the end of this chapter:} \emph{I can talk about sharing a flat and a building.}

Complete the last lesson of chapter 532: I can talk about sharing a flat and a building.
\end{chapteropening}

\section[die Wohngemeinschaft]{die Wohngemeinschaft — shared flat}

\label{lesson:GE-C532-wohngemeinschaft}



\begin{quote}
Before the new one: say the German for to wash up, then the German for to dry.
\end{quote}

\subsection*{You'll want to know: die Wohngemeinschaft}
\textbf{die Wohngemeinschaft} — \textquotedblleft{}shared flat\textquotedblright{}.

\begin{quote}
Während des Studiums wohnte ich in einer Wohngemeinschaft.
\end{quote}

\begin{quote}
\emph{While I was a student I lived in a shared flat.}
\end{quote}

Plural: \emph{die Wohngemeinschaften}. Mostly shortened to \emph{die WG}.

\begin{cousinweb}
\emph{Wohngemeinschaft} is built on \emph{wohnen} \textquotedblleft{}to live\textquotedblright{}, inherited from Old High German \emph{wonēn}, and \emph{Gemeinschaft} \textquotedblleft{}community\textquotedblright{}, from \emph{gemein} \textquotedblleft{}common\textquotedblright{} (EtymDB).
\end{cousinweb}

\begin{grammarlens}[title={What you've built}]
One more word for telling, explaining and describing at B1.
\end{grammarlens}

\subsection*{Your turn}
\begin{itemize}
\raggedright
  \item \emph{Say it:} \emph{die Wohngemeinschaft}
  \item \emph{Say it:} \emph{die Wohngemeinschaft}, once more
  \item \emph{Say it:} a sentence of your own with \emph{die Wohngemeinschaft}
  \item \emph{Recall:} say the German for to wash up, then the German for to dry, then say \emph{die Wohngemeinschaft} again
\end{itemize}

\begin{tcolorbox}[breakable,colback=teal!4,colframe=teal!35!black,title={Before you move on}]
What does die Wohngemeinschaft mean? (\textquotedblleft{}Shared flat\textquotedblright{}.) Say it once more.
\end{tcolorbox}
\section[der Mitbewohner]{der Mitbewohner — flatmate}

\label{lesson:GE-C532-mitbewohner}



\begin{quote}
Before the new one: say the German for to dry, then the German for shared flat.
\end{quote}

\subsection*{You'll want to know: der Mitbewohner}
\textbf{der Mitbewohner} — \textquotedblleft{}flatmate\textquotedblright{}.

\begin{quote}
Mein Mitbewohner kocht jeden Abend Nudeln.
\end{quote}

\begin{quote}
\emph{My flatmate cooks pasta every evening.}
\end{quote}

Plural: \emph{die Mitbewohner}. A woman is \emph{die Mitbewohnerin}.

\begin{cousinweb}
\emph{Mitbewohner} is built on \emph{mit} \textquotedblleft{}with\textquotedblright{} and \emph{Bewohner} \textquotedblleft{}resident\textquotedblright{}, from \emph{bewohnen} \textquotedblleft{}to live in\textquotedblright{} (EtymDB), built on \emph{wohnen}.
\end{cousinweb}

\begin{grammarlens}[title={What you've built}]
One more word for telling, explaining and describing at B1.
\end{grammarlens}

\subsection*{Your turn}
\begin{itemize}
\raggedright
  \item \emph{Say it:} \emph{der Mitbewohner}
  \item \emph{Say it:} \emph{der Mitbewohner}, once more
  \item \emph{Say it:} a sentence of your own with \emph{der Mitbewohner}
  \item \emph{Recall:} say the German for to dry, then the German for shared flat, then say \emph{der Mitbewohner} again
\end{itemize}

\begin{tcolorbox}[breakable,colback=teal!4,colframe=teal!35!black,title={Before you move on}]
What does der Mitbewohner mean? (\textquotedblleft{}Flatmate\textquotedblright{}.) Say it once more.
\end{tcolorbox}
\section[die Hausordnung]{die Hausordnung — house rules}

\label{lesson:GE-C532-hausordnung}



\begin{quote}
Before the new one: say the German for shared flat, then the German for flatmate.
\end{quote}

\subsection*{You'll want to know: die Hausordnung}
\textbf{die Hausordnung} — \textquotedblleft{}house rules\textquotedblright{}.

\begin{quote}
Nach der Hausordnung muss es ab zweiundzwanzig Uhr ruhig sein.
\end{quote}

\begin{quote}
\emph{Under the house rules it has to be quiet from ten p.m.}
\end{quote}

Plural: \emph{die Hausordnungen}.

\begin{cousinweb}
\emph{Hausordnung} is built on \emph{Haus} \textquotedblleft{}house\textquotedblright{} and \emph{Ordnung} \textquotedblleft{}order\textquotedblright{}, inherited from Old High German \emph{ordinunga} (EtymDB), built on a verb borrowed from Latin \emph{ōrdināre}.
\end{cousinweb}

\begin{grammarlens}[title={What you've built}]
One more word for telling, explaining and describing at B1.
\end{grammarlens}

\subsection*{Your turn}
\begin{itemize}
\raggedright
  \item \emph{Say it:} \emph{die Hausordnung}
  \item \emph{Say it:} \emph{die Hausordnung}, once more
  \item \emph{Say it:} a sentence of your own with \emph{die Hausordnung}
  \item \emph{Recall:} say the German for shared flat, then the German for flatmate, then say \emph{die Hausordnung} again
\end{itemize}

\begin{tcolorbox}[breakable,colback=teal!4,colframe=teal!35!black,title={Before you move on}]
What does die Hausordnung mean? (\textquotedblleft{}House rules\textquotedblright{}.) Say it once more.
\end{tcolorbox}
\section[der Hausmeister]{der Hausmeister — caretaker, janitor}

\label{lesson:GE-C532-hausmeister}



\begin{quote}
Before the new one: say the German for flatmate, then the German for house rules.
\end{quote}

\subsection*{You'll want to know: der Hausmeister}
\textbf{der Hausmeister} — \textquotedblleft{}caretaker, janitor\textquotedblright{}.

\begin{quote}
Der Hausmeister repariert die Tür im Keller.
\end{quote}

\begin{quote}
\emph{The caretaker is repairing the door in the cellar.}
\end{quote}

Plural: \emph{die Hausmeister}. In Switzerland: \emph{der Hauswart}.

\begin{cousinweb}
\emph{Hausmeister} is built on \emph{Haus} \textquotedblleft{}house\textquotedblright{} and \emph{Meister} \textquotedblleft{}master\textquotedblright{}, which came through Old High German \emph{meistar} from Latin \emph{magister} (EtymDB).
\end{cousinweb}

\begin{grammarlens}[title={What you've built}]
One more word for telling, explaining and describing at B1.
\end{grammarlens}

\subsection*{Your turn}
\begin{itemize}
\raggedright
  \item \emph{Say it:} \emph{der Hausmeister}
  \item \emph{Say it:} \emph{der Hausmeister}, once more
  \item \emph{Say it:} a sentence of your own with \emph{der Hausmeister}
  \item \emph{Recall:} say the German for flatmate, then the German for house rules, then say \emph{der Hausmeister} again
\end{itemize}

\begin{tcolorbox}[breakable,colback=teal!4,colframe=teal!35!black,title={Before you move on}]
What does der Hausmeister mean? (\textquotedblleft{}Caretaker, janitor\textquotedblright{}.) Say it once more.
\end{tcolorbox}
\section[das Treppenhaus]{das Treppenhaus — stairwell}

\label{lesson:GE-C532-treppenhaus}



\begin{quote}
Before the new one: say the German for house rules, then the German for caretaker.
\end{quote}

\subsection*{You'll want to know: das Treppenhaus}
\textbf{das Treppenhaus} — \textquotedblleft{}stairwell\textquotedblright{}.

\begin{quote}
Im Treppenhaus riecht es nach Essen.
\end{quote}

\begin{quote}
\emph{The stairwell smells of food.}
\end{quote}

Plural: \emph{die Treppenhäuser}.

\begin{cousinweb}
\emph{Treppenhaus} is built on \emph{Treppe} \textquotedblleft{}stairs\textquotedblright{}, usually traced to Low German, and \emph{Haus} \textquotedblleft{}house\textquotedblright{}; EtymDB lists \emph{Treppe} without a source.
\end{cousinweb}

\begin{grammarlens}[title={What you've built}]
One more word for telling, explaining and describing at B1.
\end{grammarlens}

\subsection*{Your turn}
\begin{itemize}
\raggedright
  \item \emph{Say it:} \emph{das Treppenhaus}
  \item \emph{Say it:} \emph{das Treppenhaus}, once more
  \item \emph{Say it:} a sentence of your own with \emph{das Treppenhaus}
  \item \emph{Recall:} say the German for house rules, then the German for caretaker, then say \emph{das Treppenhaus} again
\end{itemize}

\begin{tcolorbox}[breakable,colback=teal!4,colframe=teal!35!black,title={Before you move on}]
What does das Treppenhaus mean? (\textquotedblleft{}Stairwell\textquotedblright{}.) Say it once more.
\end{tcolorbox}
